% Protocolo de control TeleCon — referencia técnica + ejemplos Bluetooth.
% Compilar: pdflatex -output-directory=../build telecon_control_protocol.tex
\documentclass[12pt,a4paper]{article}

\usepackage[margin=2.4cm]{geometry}
\usepackage[T1]{fontenc}
\usepackage[utf8]{inputenc}
\usepackage{lmodern}
\usepackage[spanish]{babel}
\usepackage{booktabs}
\usepackage{hyperref}
\usepackage{xcolor}
\usepackage{colortbl}
\usepackage{array}
\usepackage{enumitem}
\usepackage{needspace}
\usepackage{etoolbox}
\usepackage{tikz}
\usepackage{listings}
\usepackage{verbatim}

\usetikzlibrary{arrows.meta,calc,fit}
\preto\section{\Needspace{10\baselineskip}}

\hypersetup{
  colorlinks=true,
  linkcolor=black,
  urlcolor=blue,
  pdftitle={Protocolo de control TeleCon},
  pdfauthor={TeleCon}
}

% Colores de sintaxis Arduino IDE 2 (legibles sobre el fondo claro del listado).
\definecolor{ArduinoKeyword}{HTML}{00979C}
\definecolor{ArduinoFunc}{HTML}{D35400}
\definecolor{ArduinoString}{HTML}{D35400}
\definecolor{ArduinoComment}{HTML}{5D6D6E}
\definecolor{ArduinoPreproc}{HTML}{9B59B6}

% Colores semánticos para tablas de campos / tramas.
\definecolor{TblHead}{HTML}{0F6F73}
\definecolor{TblSync}{HTML}{E8DAEF}
\definecolor{TblChan}{HTML}{D5F5E3}
\definecolor{TblKnob}{HTML}{D6EAF8}
\definecolor{TblSw}{HTML}{FCF3CF}
\definecolor{TblSum}{HTML}{FAD7A0}
\definecolor{TblEnd}{HTML}{E5E7E9}
\definecolor{TblTelem}{HTML}{D4E6F1}
\definecolor{TblId}{HTML}{F5CBA7}
\definecolor{TblLen}{HTML}{FDEBD0}

\newcommand{\thc}[1]{\textcolor{white}{\bfseries #1}}
\newcommand{\chip}[2]{\colorbox{#1}{\strut\,\scriptsize #2\,}}
\newcommand{\mapcell}[3]{%
  \cellcolor{#1}%
  \begin{tabular}{@{}c@{}}
    \strut\textbf{\texttt{#2}}\\[-1pt]
    \scriptsize #3
  \end{tabular}%
}

\definecolor{WireApp}{HTML}{0F6F73}
\definecolor{WireType}{HTML}{9B59B6}
\definecolor{WireKey}{HTML}{1A5276}
\definecolor{WireVal}{HTML}{D35400}
\definecolor{WireSep}{HTML}{7F8C8D}
\definecolor{WireBg}{HTML}{F4FBFC}

\newcommand{\wapp}[1]{\textcolor{WireApp}{\textbf{\texttt{#1}}}}
\newcommand{\wtyp}[1]{\textcolor{WireType}{\textbf{\texttt{#1}}}}
\newcommand{\wkey}[1]{\textcolor{WireKey}{\texttt{#1}}}
\newcommand{\wval}[1]{\textcolor{WireVal}{\texttt{#1}}}
\newcommand{\wsep}{\textcolor{WireSep}{\texttt{,}}\hspace{0pt}}
\newcommand{\wcol}{\textcolor{WireSep}{\texttt{:}}}

\tikzset{
  telecon actor/.style={
    rounded corners=4pt,
    minimum width=2.7cm,
    minimum height=1.05cm,
    font=\sffamily\bfseries,
    text=white,
    inner sep=4pt
  },
  telecon step/.style={
    circle,
    inner sep=1pt,
    minimum size=0.52cm,
    font=\sffamily\bfseries\scriptsize,
    text=white
  },
  telecon pill/.style={
    rounded corners=2.5pt,
    inner sep=3.2pt,
    font=\ttfamily\scriptsize,
    align=center
  },
  telecon arr/.style={-{Latex[length=2.1mm]}, line width=0.9pt}
}

\newcommand{\wireframe}[1]{%
  \par\addvspace{0.55em}%
  \noindent
  {\setlength{\fboxrule}{1.1pt}%
   \setlength{\fboxsep}{8pt}%
   \fcolorbox{WireApp}{WireBg}{%
     \begin{minipage}{\dimexpr\linewidth-16pt-2.2pt\relax}
     \raggedright #1
     \end{minipage}}}%
  \par\addvspace{0.35em}%
}

\lstdefinestyle{Arduino}{
  language=C++,
  basicstyle=\ttfamily\normalsize\color{black},
  identifierstyle=\ttfamily\normalsize\color{black},
  keywordstyle=\ttfamily\normalsize\color{ArduinoKeyword}\bfseries,
  keywordstyle=[2]\ttfamily\normalsize\color{ArduinoFunc},
  keywordstyle=[3]\ttfamily\normalsize\color{ArduinoKeyword},
  commentstyle=\ttfamily\normalsize\color{ArduinoComment}\itshape,
  stringstyle=\ttfamily\normalsize\color{ArduinoString},
  directivestyle=\ttfamily\normalsize\color{ArduinoPreproc}\bfseries,
  morekeywords={
    String,boolean,byte,word,nullptr,size_t,
    uint8_t,uint16_t,uint32_t,uint64_t,
    int8_t,int16_t,int32_t,int64_t
  },
  morekeywords=[2]{
    setup,loop,print,println,begin,available,read,trim,
    indexOf,substring,length,startsWith,toInt,c_str,hasClient,
    strtol,valueOf,sendLine,applyOutputs,handleConnect,handleCtrl,
    handleLine,sendTelemetry,sendPlot,sendData,toPlotByte,
    millis,constrain,map,abs,snprintf,analogRead
  },
    morekeywords=[3]{SERIAL_8N1,SERIAL_BAUD,HIGH,LOW,INPUT,OUTPUT,INPUT_PULLUP}
}

\lstset{
  style=Arduino,
  breaklines=true,
  columns=fullflexible,
  frame=single,
  xleftmargin=0.3em,
  xrightmargin=0.3em,
  framesep=3pt,
  backgroundcolor=\color{gray!8},
  showstringspaces=false,
  keepspaces=true
}

\title{Protocolo de control TeleCon\\[0.4em]
\large Contrato de flujo de bytes con la aplicación Android}
\author{Referencia de firmware --- cualquier MCU que implemente este contrato}
\date{Versión 1.2 \quad 2026}

\begin{document}
\maketitle

\begin{abstract}
Este documento especifica cómo el firmware debe intercambiar bytes con
TeleCon4ESP32 por Bluetooth Classic SPP, BLE Nordic UART o TCP.
La aplicación es agnóstica al transporte: no exige un SoC Espressif, un
nombre GAP concreto ni comandos AT de fabricante.
Los sketches oficiales de Arduino son implementaciones de referencia de
este contrato, no parte de la definición en el cable.

Control Panel y RC Vehicle Pro comparten el prefijo de aplicación \texttt{RC}.
El vídeo HTTP de cámara (\texttt{/stream}, \texttt{/capture}) es una sesión
aparte y queda fuera de alcance.
\end{abstract}

\section{Arquitectura}
\label{sec:arch}

El teléfono es siempre el cliente. Tras establecer el enlace subyacente,
abre una sesión de aplicación. En el dispositivo hay que implementar tres
capas:

\begin{enumerate}[leftmargin=1.6cm]
  \item \textbf{Transporte} --- un flujo de octetos bidireccional
        suficientemente fiable (sección~\ref{sec:pipes}).
  \item \textbf{Handshake de sesión} --- una línea de texto \texttt{CONNECT}
        y una respuesta \texttt{ACK}/\texttt{NAK} (sección~\ref{sec:handshake}).
        El tiempo de espera en el teléfono es $2{,}5$\,s. Hasta el ACK, el
        control de la UI no es una sesión TeleCon viva, aunque el SO marque
        el socket como conectado.
  \item \textbf{Carga} --- líneas de texto Simple (sección~\ref{sec:simple})
        o tramas Binary (sección~\ref{sec:binary}), según el ajuste del
        usuario y el campo \texttt{proto} de \texttt{CONNECT}.
\end{enumerate}

\Needspace{13.2cm}
\begin{center}
\begin{tikzpicture}[x=1cm, y=0.78cm]
  \fill[rounded corners=7pt, WireBg] (-0.4,-1.45) rectangle (14.15,8.7);
  \draw[rounded corners=7pt, TblHead!28, line width=0.7pt]
    (-0.4,-1.45) rectangle (14.15,8.7);

  \fill[TblHead!11, rounded corners=5pt] (0.08,5.62) rectangle (13.72,7.02);
  \fill[ArduinoPreproc!10, rounded corners=5pt] (0.08,3.62) rectangle (13.72,5.52);
  \fill[ArduinoKeyword!10, rounded corners=5pt] (0.08,1.55) rectangle (13.72,3.52);
  \fill[ArduinoFunc!11, rounded corners=5pt] (0.08,-0.12) rectangle (13.72,1.45);

  \node[telecon actor, fill=TblHead, align=center] at (2.4,7.85)
    {Tel\'efono\\[-1.6pt]{\normalfont\sffamily\scriptsize cliente}};
  \node[telecon actor, fill=ArduinoPreproc, align=center] at (11.8,7.85)
    {MCU\\[-1.6pt]{\normalfont\sffamily\scriptsize firmware}};
  \draw[TblHead, line width=1.4pt] (2.4,7.22) -- (2.4,0.12);
  \draw[ArduinoPreproc, line width=1.4pt] (11.8,7.22) -- (11.8,0.12);

  \node[telecon step, fill=TblHead] at (0.55,6.42) {1};
  \draw[telecon arr, TblHead] (2.6,6.55) -- (11.6,6.55);
  \node[telecon pill, fill=white, draw=TblHead!55] at (7.1,6.55)
    {SPP \textbullet\ NUS \textbullet\ TCP};
  \node[font=\sffamily\scriptsize\bfseries, text=TblHead] at (7.1,5.88)
    {1. Abrir el enlace};

  \node[telecon step, fill=ArduinoPreproc] at (0.55,4.65) {2};
  \draw[telecon arr, ArduinoPreproc] (2.6,5.12) -- (11.6,5.12);
  \node[telecon pill, fill=TblSync] at (7.1,5.12) {RC:CONNECT};
  \draw[telecon arr, ArduinoKeyword] (11.6,4.28) -- (2.6,4.28);
  \node[telecon pill, fill=TblChan] at (5.25,4.28) {RC:ACK};
  \node[font=\sffamily\scriptsize\bfseries, text=ArduinoComment] at (7.1,4.28)
    {o};
  \node[telecon pill, fill=TblId] at (8.95,4.28) {RC:NAK};
  \node[font=\sffamily\scriptsize\bfseries, text=ArduinoPreproc] at (7.1,3.82)
    {2. Handshake ($2{,}5$\,s) --- el ACK abre la sesi\'on};

  \node[telecon step, fill=ArduinoKeyword] at (0.55,2.6) {3};
  \draw[telecon arr, ArduinoKeyword] (2.6,3.12) -- (11.6,3.12);
  \node[telecon pill, fill=TblChan] at (7.1,3.12) {CTRL / AA\,55 / BTN};
  \draw[telecon arr, ArduinoFunc] (11.6,2.22) -- (2.6,2.22);
  \node[telecon pill, fill=TblTelem] at (7.1,2.22) {DATA / PLOT / CC\,11--33};
  \node[font=\sffamily\scriptsize\bfseries, text=ArduinoKeyword] at (7.1,1.75)
    {3. Tras el ACK: control hacia abajo, telemetr\'ia hacia arriba};

  \node[telecon step, fill=ArduinoFunc] at (0.55,0.72) {4};
  \draw[telecon arr, ArduinoFunc] (2.6,0.85) -- (11.6,0.85);
  \node[telecon pill, fill=TblSum] at (7.1,0.85) {cae el enlace};
  \node[font=\sffamily\scriptsize\bfseries, text=ArduinoFunc] at (7.1,0.18)
    {4. Fail-safe en el MCU};

  \node[font=\sffamily\scriptsize, text=ArduinoComment] at (7.0,-0.95)
    {El tel\'efono siempre empieza. Hasta el ACK, el movimiento de la UI no es sesi\'on viva.};
\end{tikzpicture}

\vspace{0.3em}
{\sffamily\scriptsize
\chip{TblHead!18}{\textcolor{TblHead}{\textbf{1} enlace}}
\chip{TblSync}{\textbf{2} CONNECT}
\chip{TblChan}{ACK / CTRL}
\chip{TblId}{NAK}
\chip{TblTelem}{DATA / PLOT}
\chip{TblSum}{\textbf{4} ca\'ida}}

\vspace{0.15em}
{\sffamily\scriptsize\itshape
Flujo tel\'efono--MCU. Verde = aceptar, naranja = rechazar, verde azulado = control, azul = telemetr\'ia.}
\end{center}

El nombre anunciado por Bluetooth (por ejemplo
\texttt{ESP32-TC-RC-BT-Simple} en el firmware oficial) es solo presentación
GAP/SDP.
El analizador del handshake no busca la cadena \texttt{ESP32}.
La conexión usa la \textbf{dirección} remota más el UUID SPP o los
atributos NUS.

\subsection{Acoplamiento de modos}

El transporte y la carga del firmware deben coincidir con el
\texttt{BluetoothConnectionMode} almacenado:

\begin{center}
\renewcommand{\arraystretch}{1.25}
\begin{tabular}{>{\raggedright\arraybackslash}p{3.5cm}
                >{\centering\arraybackslash}p{2.3cm}
                >{\centering\arraybackslash}p{1.8cm}
                >{\raggedright\arraybackslash}p{3.8cm}}
\rowcolor{TblHead}
\thc{Modo en la aplicación} & \thc{Transporte} & \thc{proto} & \thc{Carga tras ACK} \\
\rowcolor{TblChan}
Classic Simple & SPP & \texttt{simple} & Sección~\ref{sec:simple} \\
\rowcolor{TblTelem}
Classic Binary & SPP & \texttt{binary} & Sección~\ref{sec:binary} \\
\rowcolor{TblTelem}
BLE Binary & NUS & \texttt{binary} & Sección~\ref{sec:binary} \\
\rowcolor{TblTelem}
Wi-Fi SoftAP Binary & TCP :3333 & \texttt{binary} & Sección~\ref{sec:binary} \\
\rowcolor{TblChan}
CAM starter (rol~B) & TCP :3333 & \texttt{simple} & Sección~\ref{sec:simple} \\
\rowcolor{TblTelem}
CAM SoftAP Binary & TCP :3333 & \texttt{binary} & Sección~\ref{sec:binary} \\
\rowcolor{TblId}
CAM SoftAP heredado & TCP :3333 & \texttt{wifi} & Texto Simple \\
\end{tabular}

\vspace{0.4em}
\chip{TblChan}{Simple}
\chip{TblTelem}{Binary}
\chip{TblId}{Simple heredado}
\end{center}

BLE es solo Binary en la aplicación actual.
Un perfil GATT serie propio (UUID típicos de HM-10) falla en el descubrimiento:
el cliente solo busca Nordic UART.

La discrepancia debe señalarse con \texttt{NAK} (\texttt{expected} = capacidad
del dispositivo, \texttt{actual} = valor de \texttt{CONNECT}).
Ignorarla produce \texttt{HandshakeFailure.Timeout} en el teléfono.

\section{Transportes}
\label{sec:pipes}

\subsection{Classic SPP (RFCOMM)}

Serial Port Profile, UUID
\texttt{00001101-0000-1000-8000-00805F9B34FB}.

El cliente Android intenta, en este orden:
\texttt{createRfcommSocketToServiceRecord} seguro,
\texttt{createInsecureRfcommSocketToServiceRecord} inseguro y luego un
socket RFCOMM por reflexión en el canal~1.
El tercer intento existe porque algunos stacks (incluido
\texttt{BluetoothSerial} del ESP32 con un \texttt{WifiNetworkSpecifier}
SoftAP activo) fallan la búsqueda SDP del UUID.

HC-05 / HC-06 exponen este UUID como puente UART transparente.
El MCU solo implementa el protocolo de bytes en ese UART; el baud es
ajuste del módulo (habitualmente 9600 o 115200), no lo negocia el teléfono.

\subsection{BLE Nordic UART Service}

\begin{center}
\renewcommand{\arraystretch}{1.25}
\begin{tabular}{>{\raggedright\arraybackslash}p{4.4cm}
                >{\raggedright\arraybackslash}p{6.2cm}
                >{\centering\arraybackslash}p{2.7cm}}
\rowcolor{TblHead}
\thc{Atributo en el periférico} & \thc{UUID} & \thc{Rol} \\
\rowcolor{TblSync}
Servicio primario &
  {\small\ttfamily 6E400001-\allowbreak B5A3-F393-\allowbreak E0A9-\allowbreak E50E24DCCA9E} &
  Servicio \\
\rowcolor{TblChan}
RX (Write / Write Without Response) &
  {\small\ttfamily 6E400002-\ldots CCA9E} &
  Teléfono $\rightarrow$ dispositivo \\
\rowcolor{TblTelem}
TX (Notify) &
  {\small\ttfamily 6E400003-\ldots CCA9E} &
  Dispositivo $\rightarrow$ teléfono \\
\rowcolor{TblLen}
CCCD en TX &
  {\small\ttfamily 00002902-\allowbreak 0000-1000-\allowbreak 8000-\allowbreak 00805F9B34FB} &
  Activar notify \\
\end{tabular}

\vspace{0.4em}
\chip{TblSync}{Servicio}
\chip{TblChan}{El teléfono escribe}
\chip{TblTelem}{El dispositivo notifica}
\chip{TblLen}{Config. cliente}
\end{center}

El cliente escribe octetos de protocolo en RX y recompone las notificaciones
TX en el mismo flujo que SPP.
Cada notify/write debe ser $\le$ MTU$-3$.
El MTU ATT por defecto (23) implica trozos de 20 bytes; pida un MTU mayor
si envía tramas Binary de una sola vez.
El handshake sigue siendo ASCII delimitado por newline en este flujo.

\subsection{TCP}

La misma sesión y carga que Classic, en una conexión TCP.
El firmware SoftAP oficial enlaza \texttt{192.168.4.1:3333}.
Host y puerto no forman parte del protocolo de aplicación; son parámetros
de despliegue.
No invente un segundo handshake para TCP.

\section{Codificación de texto}

Handshake y mensajes Simple. El color indica la función de cada token:

\begin{center}
\setlength{\tabcolsep}{4pt}
\renewcommand{\arraystretch}{1.15}
\begin{tabular}{|*{6}{>{\centering\arraybackslash}p{1.85cm}|}}
\hline
\mapcell{TblChan}{APP}{app} &
\mapcell{TblSync}{TYPE}{tipo} &
\mapcell{TblKnob}{key1}{clave} &
\mapcell{TblId}{value1}{valor} &
\mapcell{TblKnob}{key2}{clave} &
\mapcell{TblId}{value2}{valor} \\
\hline
\end{tabular}
\end{center}

\wireframe{%
\wapp{APP}\wcol\wtyp{TYPE}\wsep
\wkey{key1}\wsep\wval{value1}\wsep
\wkey{key2}\wsep\wval{value2}\wsep\texttt{...}%
}

\begin{center}
\vspace{-0.15em}
\chip{TblChan}{App}
\chip{TblSync}{Tipo}
\chip{TblKnob}{Clave}
\chip{TblId}{Valor}
\end{center}

Terminados en \texttt{LF} (\texttt{0x0A}). Un \texttt{CR} previo se elimina.
\texttt{APP} en este producto es \texttt{RC}.
Las claves respetan mayúsculas tal como las envía el teléfono
(\texttt{proto}, \texttt{lx}, \ldots).

Booleanos: \texttt{1}/\texttt{0}, \texttt{true}/\texttt{false},
\texttt{on}/\texttt{off} o cualquier entero distinto de cero.

El ensamblador de entrada trata un \texttt{0xCC} inicial como inicio de
trama Binary de telemetría y el ASCII imprimible como texto.
Tras el ACK en modo Binary, \texttt{0xAA}/\texttt{0xBB} son tramas de
control hacia el dispositivo.

\section{Handshake de sesión}
\label{sec:handshake}

Obligatorio en cada conexión nueva SPP, NUS o TCP.
No se usa para la cámara HTTP.

\subsection{Cliente $\rightarrow$ dispositivo}

\wireframe{%
\wapp{RC}\wcol\wtyp{CONNECT}\wsep\wkey{proto}\wsep\wval{simple}\\[0.25em]
\wapp{RC}\wcol\wtyp{CONNECT}\wsep\wkey{proto}\wsep\wval{binary}\\[0.25em]
\wapp{RC}\wcol\wtyp{CONNECT}\wsep\wkey{proto}\wsep\wval{wifi}%
}

\begin{center}
\vspace{-0.15em}
\chip{TblChan}{App}
\chip{TblSync}{Tipo}
\chip{TblKnob}{Clave}
\chip{TblId}{Valor}
\end{center}

\texttt{wifi} se conserva para SoftAP CAM de una sola placa heredado.
El firmware Simple nuevo debe aceptar \texttt{simple} y \texttt{wifi} y
emitir el mismo ACK.

El teléfono espera $2500$\,ms una línea con prefijo \texttt{RC} y tipo
\texttt{ACK} o \texttt{NAK}.

\subsection{Dispositivo $\rightarrow$ cliente}

\wireframe{%
\wapp{RC}\wcol\wtyp{ACK}\wsep\wkey{app}\wsep\wval{RC}%
}

Las claves extra del ACK se ignoran al establecer la sesión.

\wireframe{%
\wapp{RC}\wcol\wtyp{NAK}\wsep
\wkey{reason}\wsep\wval{proto\_mismatch}\wsep
\wkey{expected}\wsep\wval{simple}\wsep
\wkey{actual}\wsep\wval{binary}\\[0.25em]
\wapp{RC}\wcol\wtyp{NAK}\wsep
\wkey{reason}\wsep\wval{app\_mismatch}\wsep
\wkey{expected}\wsep\wval{RC}\wsep
\wkey{actual}\wsep\wval{XX}%
}

\texttt{reason} debe ser \texttt{proto\_mismatch} o \texttt{app\_mismatch}
para los diálogos estructurados de Android.
\texttt{RC:ACK,steer\_center,1} es un acuse posterior de ajustes y no debe
usarse como éxito de CONNECT.

\subsection{Estado recomendado en el dispositivo}

\begin{enumerate}[leftmargin=1.6cm]
  \item Enlace arriba $\rightarrow$ \texttt{session\_ok = false}, salidas
        en fail-safe.
  \item Línea \texttt{CONNECT} completa $\rightarrow$ comparar \texttt{proto}
        (y, si se desea, el prefijo de app) $\rightarrow$ ACK o NAK.
  \item Tras ACK, habilitar \texttt{CTRL}/\texttt{AA 55} y TX de telemetría.
  \item Al caer el enlace, resetear \texttt{session\_ok} y las salidas.
\end{enumerate}

\section{\texorpdfstring{Implementación de referencia: SPP Simple\\[0.15em](ESP32 \texttt{BluetoothSerial})}{Implementación de referencia: SPP Simple (ESP32 BluetoothSerial)}}
\label{sec:code-esp32}

Bluetooth Classic en piezas clase ESP32-WROOM.
Ajuste de la aplicación: Classic Simple.
El nombre GAP es arbitrario.

\begin{lstlisting}
#include "BluetoothSerial.h"

#define SERIAL_BAUD         115200

BluetoothSerial SerialBT;

const char* BT_NAME = "TeleCon-RC";
const char* MY_PROTO = "simple";

bool sessionOk = false;
String lineBuf;

int lx = 0, ly = 0, rx = 0, ry = 0;
int lk = 512, rk = 512;
uint8_t sw = 0;

String valueOf(const String& line, const char* key) {
  String token = String(key) + ",";
  int i = line.indexOf(token);
  if (i < 0) return "";
  i += token.length();
  int j = line.indexOf(',', i);
  if (j < 0) j = line.length();
  return line.substring(i, j);
}

void sendLine(const String& s) {
  SerialBT.print(s);
  SerialBT.print('\n');
}

void applyOutputs() {
  // Mapear palancas ya condicionadas (-100..100) a actuadores.
}

void handleConnect(const String& line) {
  String proto = valueOf(line, "proto");
  if (proto != MY_PROTO &&
      !(String(MY_PROTO) == "simple" && proto == "wifi")) {
    sendLine(String("RC:NAK,reason,proto_mismatch,expected,") +
             MY_PROTO + ",actual," + proto);
    sessionOk = false;
    return;
  }
  sendLine("RC:ACK,app,RC");
  sessionOk = true;
}

void handleCtrl(const String& line) {
  lx = valueOf(line, "lx").toInt();
  ly = valueOf(line, "ly").toInt();
  rx = valueOf(line, "rx").toInt();
  ry = valueOf(line, "ry").toInt();
  lk = valueOf(line, "lk").toInt();
  rk = valueOf(line, "rk").toInt();
  sw = (uint8_t)strtol(valueOf(line, "sw").c_str(), nullptr, 16);
  applyOutputs();
}

void handleLine(String line) {
  line.trim();
  if (line.length() == 0) return;
  if (line.startsWith("RC:CONNECT")) {
    handleConnect(line);
    return;
  }
  if (!sessionOk) return;
  if (line.startsWith("RC:CTRL")) handleCtrl(line);
  else if (line.startsWith("RC:BTN")) {
    (void)valueOf(line, "id").toInt();
  } else if (line.startsWith("RC:SET")) {
    // persistir steer_center / rx si implementa trim
  }
}

void setup() {
  Serial.begin(SERIAL_BAUD);
  SerialBT.begin(BT_NAME);
}

void loop() {
  if (!SerialBT.hasClient()) {
    if (sessionOk) {
      sessionOk = false;
      // fail-safe
    }
    lineBuf = "";
  }
  while (SerialBT.available()) {
    char c = (char)SerialBT.read();
    if (c == '\n') {
      handleLine(lineBuf);
      lineBuf = "";
    } else if (c != '\r') {
      if (lineBuf.length() < 200) lineBuf += c;
      else lineBuf = "";
    }
  }
}
\end{lstlisting}

Telemetría tras el ACK (DATA típico $2$\,Hz, PLOT $10$--$20$\,Hz):

\begin{lstlisting}
void sendTelemetry(int leftPanel, int batt, int analog) {
  sendLine("RC:DATA,left," + String(leftPanel) +
           ",batt," + String(batt) + ",analog," + String(analog));
}

void sendPlot(uint8_t v0, uint8_t v1, uint8_t v2, uint8_t v3) {
  sendLine("RC:PLOT,v0," + String(v0) + ",v1," + String(v1) +
           ",v2," + String(v2) + ",v3," + String(v3));
}
\end{lstlisting}

\section{\texorpdfstring{Implementación de referencia: SPP Simple\\[0.15em]por UART HC-05}{Implementación de referencia: SPP Simple por UART HC-05}}
\label{sec:code-hc05}

El módulo termina SPP; el MCU ve un UART.
La misma gramática \texttt{CONNECT}/\texttt{CTRL}.
Iguale el baud del módulo (a menudo 9600). Cruce TX/RX; adapte niveles
si hace falta.

\begin{lstlisting}
#include <Arduino.h>

#define SERIAL_BAUD         115200

HardwareSerial BT(1);
const int BT_RX_PIN = 16;
const int BT_TX_PIN = 17;
const uint32_t BT_BAUD = 9600;
const char* MY_PROTO = "simple";

bool sessionOk = false;
String lineBuf;

String valueOf(const String& line, const char* key) {
  String token = String(key) + ",";
  int i = line.indexOf(token);
  if (i < 0) return "";
  i += token.length();
  int j = line.indexOf(',', i);
  if (j < 0) j = line.length();
  return line.substring(i, j);
}

void sendLine(const String& s) {
  BT.print(s);
  BT.print('\n');
}

void handleConnect(const String& line) {
  String proto = valueOf(line, "proto");
  if (proto != MY_PROTO) {
    sendLine(String("RC:NAK,reason,proto_mismatch,expected,") +
             MY_PROTO + ",actual," + proto);
    sessionOk = false;
    return;
  }
  sendLine("RC:ACK,app,RC");
  sessionOk = true;
}

void handleLine(String line) {
  line.trim();
  if (line.startsWith("RC:CONNECT")) { handleConnect(line); return; }
  if (!sessionOk) return;
  if (line.startsWith("RC:CTRL")) {
    int ly = valueOf(line, "ly").toInt();
    int rxStick = valueOf(line, "rx").toInt();
    (void)ly;
    (void)rxStick;
  }
}

void setup() {
  Serial.begin(SERIAL_BAUD);
  BT.begin(BT_BAUD, SERIAL_8N1, BT_RX_PIN, BT_TX_PIN);
}

void loop() {
  while (BT.available()) {
    char c = (char)BT.read();
    if (c == '\n') { handleLine(lineBuf); lineBuf = ""; }
    else if (c != '\r') {
      if (lineBuf.length() < 200) lineBuf += c;
      else lineBuf = "";
    }
  }
}
\end{lstlisting}

\section{Carga Simple}
\label{sec:simple}

\subsection{Teléfono $\rightarrow$ dispositivo}

\paragraph{\texttt{RC:CTRL}}
Se emite al cambiar la UI (no es un sondeo fijo en el teléfono).

\wireframe{%
\wapp{RC}\wcol\wtyp{CTRL}\wsep
\wkey{lx}\wsep\wval{-50}\wsep
\wkey{ly}\wsep\wval{100}\wsep
\wkey{rx}\wsep\wval{0}\wsep
\wkey{ry}\wsep\wval{0}\wsep
\wkey{lk}\wsep\wval{512}\wsep
\wkey{rk}\wsep\wval{256}\wsep
\wkey{sw}\wsep\wval{03}%
}

\begin{center}
\renewcommand{\arraystretch}{1.25}
\begin{tabular}{>{\centering\arraybackslash}p{3.6cm}
                >{\centering\arraybackslash}p{4.2cm}
                >{\raggedright\arraybackslash}p{4.8cm}}
\rowcolor{TblHead}
\thc{Clave} & \thc{Dominio} & \thc{Significado} \\
\rowcolor{TblChan}
\texttt{lx}, \texttt{ly}, \texttt{rx}, \texttt{ry} & $-100\ldots 100$ & Palancas \\
\rowcolor{TblKnob}
\texttt{lk}, \texttt{rk} & $0\ldots 1023$ & Potenciómetros \\
\rowcolor{TblSw}
\texttt{sw} & hex 8 bits, dos dígitos & S1\,=\,bit0 \ldots\ S8\,=\,bit7 \\
\end{tabular}

\vspace{0.4em}
\chip{TblChan}{Palancas}
\chip{TblKnob}{Potenciómetros}
\chip{TblSw}{Interruptores}
\end{center}

Dual-rate, expo, inverso y zona muerta se aplican en el teléfono antes del TX.
El firmware no debe volver a aplicar esas curvas.

\paragraph{\texttt{RC:BTN}}
\texttt{RC:BTN,id,1} --- id 1--4 cluster central; $16$ ($0x10$) guardar
centro de dirección en RC Vehicle Pro.

\paragraph{\texttt{RC:SET}}
\texttt{RC:SET,steer\_center,1,rx,-3} --- persistir el cero mecánico.
\texttt{rx} incluye el trim del teléfono ($-100\ldots 100$).

\subsection{Dispositivo $\rightarrow$ teléfono}

La telemetría es opcional. Envíela solo tras el ACK, en el mismo flujo de
texto que \texttt{CTRL}. Hay dos tipos de línea: \texttt{DATA} para paneles
e indicadores lentos, y \texttt{PLOT} para un bus analógico rápido.
Si el firmware no tiene nada que mostrar, omita el tipo; el teléfono
conserva los últimos valores.

\paragraph{\texttt{RC:DATA} --- paneles, indicadores, LED (típico $2$\,Hz)}

Cada línea es una instantánea de campos enteros del HUD. Las claves pueden
ir en cualquier orden. Si falta una clave, el teléfono mantiene el valor
anterior.

\wireframe{%
\wapp{RC}\wcol\wtyp{DATA}\wsep
\wkey{left}\wsep\wval{1234}\wsep
\wkey{right}\wsep\wval{5678}\wsep
\wkey{analog}\wsep\wval{42}\wsep
\wkey{batt}\wsep\wval{88}\wsep
\wkey{led}\wsep\wval{0F}%
}

\begin{center}
\renewcommand{\arraystretch}{1.25}
\setlength{\tabcolsep}{4pt}
\begin{tabular}{>{\centering\arraybackslash}p{1.8cm}
                >{\centering\arraybackslash}p{3.2cm}
                >{\raggedright\arraybackslash}p{5.2cm}
                >{\raggedright\arraybackslash}p{3.2cm}}
\rowcolor{TblHead}
\thc{Clave} & \thc{Dominio} & \thc{Significado} & \thc{HUD típico} \\
\rowcolor{TblChan}
\texttt{left} & entero & Panel numérico izquierdo & RC Vehicle Pro: velocidad$\times 10$ \\
\rowcolor{TblChan}
\texttt{right} & entero & Panel numérico derecho & Segundo contador / sin uso \\
\rowcolor{TblTelem}
\texttt{analog} & $0\ldots 255$ & Indicador analógico & Temperatura del motor \\
\rowcolor{TblTelem}
\texttt{batt} & $0\ldots 255$ & Indicador de batería & Basta con carga $0\ldots 100$ \\
\rowcolor{TblSw}
\texttt{led} & hex 8 bits, dos dígitos & Bits de LED & Lámparas de estado \\
\end{tabular}

\vspace{0.35em}
\chip{TblChan}{Paneles}
\chip{TblTelem}{Indicadores}
\chip{TblSw}{LED}
\end{center}

El encendido y el color de los paneles se configuran en Android, no en el
cable. Las claves heredadas \texttt{lo}, \texttt{ro}, \texttt{lg},
\texttt{rg}, \texttt{lt}, \texttt{rt}, \texttt{at}, \texttt{bt} se ignoran
para la apariencia.

El teléfono puede reasignar estas claves (y los canales PLOT) a gráficas,
radar o indicadores en los ajustes RC. El firmware debe seguir enviando las
mismas claves. No emita \texttt{RC:RADAR} ni paquetes de mapa de canales.

\paragraph{\texttt{RC:PLOT} --- muestras analógicas (típico $10$--$30$\,Hz)}

Cada línea es una muestra de tiempo en el bus numerado \texttt{v0}\ldots\texttt{v7}
(= CH1\ldots CH8). Mapee sensores a $0\ldots 255$ ($0$ = abajo de la gráfica,
$255$ = arriba). Una línea de cuatro muestras es válida; si omite
\texttt{v4}\ldots\texttt{v7}, el teléfono los deja en $0$. La aplicación
guarda los últimos $100$ puntos por canal.

\wireframe{%
\wapp{RC}\wcol\wtyp{PLOT}\wsep
\wkey{v0}\wsep\wval{128}\wsep
\wkey{v1}\wsep\wval{200}\wsep
\wkey{v2}\wsep\wval{64}\wsep
\wkey{v3}\wsep\wval{180}\\[0.25em]
\wapp{RC}\wcol\wtyp{PLOT}\wsep
\wkey{v0}\wsep\wval{128}\wsep
\wkey{v1}\wsep\wval{200}\wsep
\wkey{v2}\wsep\wval{64}\wsep
\wkey{v3}\wsep\wval{180}\wsep
\wkey{v4}\wsep\wval{0}\wsep
\wkey{v5}\wsep\wval{0}\wsep
\wkey{v6}\wsep\wval{0}\wsep
\wkey{v7}\wsep\wval{0}%
}

\begin{center}
\renewcommand{\arraystretch}{1.25}
\setlength{\tabcolsep}{4pt}
\begin{tabular}{>{\centering\arraybackslash}p{2.2cm}
                >{\centering\arraybackslash}p{2.4cm}
                >{\centering\arraybackslash}p{2.2cm}
                >{\raggedright\arraybackslash}p{6.6cm}}
\rowcolor{TblHead}
\thc{Clave} & \thc{Canal} & \thc{Rango} & \thc{Notas} \\
\rowcolor{TblChan}
\texttt{v0} & CH1 & $0\ldots 255$ & Primera traza por defecto \\
\rowcolor{TblChan}
\texttt{v1} & CH2 & $0\ldots 255$ & Segunda traza / radar por defecto \\
\rowcolor{TblKnob}
\texttt{v2} & CH3 & $0\ldots 255$ & Tercera traza por defecto \\
\rowcolor{TblKnob}
\texttt{v3} & CH4 & $0\ldots 255$ & Cuarta traza por defecto \\
\rowcolor{TblTelem}
\texttt{v4}\ldots\texttt{v7} & CH5\ldots CH8 & $0\ldots 255$ & Fuentes extra; asígnelas en el mapa de canales \\
\end{tabular}

\vspace{0.35em}
\chip{TblChan}{CH1--CH2}
\chip{TblKnob}{CH3--CH4}
\chip{TblTelem}{CH5--CH8 opcionales}
\end{center}

Control Panel muestra cuatro widgets de gráfica (por defecto CH1--CH4).
Los colores de traza están fijos por índice de canal. Las etiquetas y el
mapa de canales viven en los ajustes RC de Android --- no envíe
\texttt{RC:PLOTCFG}.

\texttt{loop()} típico: PLOT cada $30$--$100$\,ms desde ADC o sensores
filtrados; DATA cada $500$\,ms para paneles y batería. Condicione ambos a
\texttt{session\_ok}.

Inserte esto en el sketch de la sección~\ref{sec:code-esp32}.
\texttt{sendLine}, \texttt{sessionOk}, \texttt{ly}, \texttt{rx} y
\texttt{sw} ya están ahí.

\begin{lstlisting}
int toPlotByte(int value, int minVal, int maxVal) {
  if (maxVal <= minVal) return 0;
  return constrain(map(value, minVal, maxVal, 0, 255), 0, 255);
}

void sendData(int left, int right, int analog, int batt, uint8_t led) {
  char line[96];
  snprintf(line, sizeof(line),
           "RC:DATA,left,%d,right,%d,analog,%d,batt,%d,led,%02X",
           left, right, analog, batt, led);
  sendLine(line);
}

void sendPlot(uint8_t v0, uint8_t v1, uint8_t v2, uint8_t v3) {
  sendLine("RC:PLOT,v0," + String(v0) + ",v1," + String(v1) +
           ",v2," + String(v2) + ",v3," + String(v3));
}

unsigned long lastDataMs = 0;
unsigned long lastPlotMs = 0;

void loop() {
  // leer SerialBT y handleLine() como en la seccion 5
  if (!sessionOk) return;

  unsigned long now = millis();
  if (now - lastDataMs >= 500) {
    lastDataMs = now;
    int speedX10 = abs(ly) * 25 / 10;  // HUD de RC Vehicle Pro
    sendData(speedX10, 0, 90, 76, sw);
  }
  if (now - lastPlotMs >= 50) {
    lastPlotMs = now;
    sendPlot(toPlotByte(analogRead(34), 0, 4095),
             toPlotByte(abs(ly), 0, 100),
             toPlotByte(76, 0, 100),
             toPlotByte(abs(rx), 0, 100));
  }
}
\end{lstlisting}

\Needspace{0.72\textheight}
\section{Carga Binary}
\label{sec:binary}

Se usa tras \texttt{RC:CONNECT,proto,binary} y ACK.
El handshake sigue siendo ASCII.

\subsection{Teléfono $\rightarrow$ dispositivo --- \texttt{AA 55} (18 bytes)}

Canales de 12 bits little-endian. El color indica la función de cada campo:

\begin{center}
\begin{minipage}{\textwidth}
\centering
\setlength{\tabcolsep}{3pt}
\renewcommand{\arraystretch}{1.15}
\begin{tabular}{|*{10}{>{\centering\arraybackslash}p{1.28cm}|}}
\hline
\mapcell{TblSync}{AA 55}{0--1} &
\mapcell{TblChan}{Lx}{2--3} &
\mapcell{TblChan}{Ly}{4--5} &
\mapcell{TblChan}{Rx}{6--7} &
\mapcell{TblChan}{Ry}{8--9} &
\mapcell{TblKnob}{Lk}{10--11} &
\mapcell{TblKnob}{Rk}{12--13} &
\mapcell{TblSw}{SW}{14} &
\mapcell{TblSum}{$\sum$}{15} &
\mapcell{TblEnd}{CR LF}{16--17} \\
\hline
\end{tabular}

\vspace{0.45em}
\renewcommand{\arraystretch}{1.25}
\setlength{\tabcolsep}{4pt}
\begin{tabular}{>{\centering\arraybackslash}p{1.35cm}
                >{\centering\arraybackslash}p{1.05cm}
                >{\raggedright\arraybackslash}p{9.15cm}
                >{\centering\arraybackslash}p{2.5cm}}
\rowcolor{TblHead}
\thc{Despl.} & \thc{Tam.} & \thc{Contenido} & \thc{Rol} \\
\rowcolor{TblSync}
0--1 & 2 & Sincronía \texttt{0xAA 0x55} & Sincronía \\
\rowcolor{TblChan}
2--3 & 2 & Lx, 12 bit LE & Canal \\
\rowcolor{TblChan}
4--5 & 2 & Ly, 12 bit LE & Canal \\
\rowcolor{TblChan}
6--7 & 2 & Rx, 12 bit LE & Canal \\
\rowcolor{TblChan}
8--9 & 2 & Ry, 12 bit LE & Canal \\
\rowcolor{TblKnob}
10--11 & 2 & Potenciómetro izquierdo $0\ldots 1023$, 12 bit LE & Potenciómetro \\
\rowcolor{TblKnob}
12--13 & 2 & Potenciómetro derecho $0\ldots 1023$, 12 bit LE & Potenciómetro \\
\rowcolor{TblSw}
14 & 1 & Bits de interruptores (S1\,=\,bit0 \ldots\ S8\,=\,bit7) & Interruptores \\
\rowcolor{TblSum}
15 & 1 & $\sum$ bytes $2\ldots 14$ (8 bits bajos) & Checksum \\
\rowcolor{TblEnd}
16--17 & 2 & Terminador \texttt{0x0D 0x0A} & Terminador \\
\end{tabular}

\vspace{0.4em}
\chip{TblSync}{Sincronía}
\chip{TblChan}{Canales}
\chip{TblKnob}{Potenciómetros}
\chip{TblSw}{Interruptores}
\chip{TblSum}{Checksum}
\chip{TblEnd}{Terminador}
\end{minipage}
\end{center}

Mapeo de palanca: porcentaje $p\in[-100,100] \mapsto \lfloor(p+100)\times 4094/200\rfloor$
con centro $0\mapsto 2047$ (unsigned $0\ldots 4094$).

\subsection{Teléfono $\rightarrow$ dispositivo --- \texttt{BB 66} (4 bytes)}

Los id coinciden con \texttt{RC:BTN} Simple.

\begin{center}
\renewcommand{\arraystretch}{1.25}
\setlength{\tabcolsep}{4pt}
\begin{tabular}{>{\centering\arraybackslash}p{1.35cm}
                >{\centering\arraybackslash}p{1.05cm}
                >{\raggedright\arraybackslash}p{9.15cm}
                >{\centering\arraybackslash}p{2.5cm}}
\rowcolor{TblHead}
\thc{Despl.} & \thc{Tam.} & \thc{Contenido} & \thc{Rol} \\
\rowcolor{TblSync}
0--1 & 2 & Sincronía \texttt{0xBB 0x66} & Sincronía \\
\rowcolor{TblId}
2 & 1 & Id de botón (1--4; $16$\,=\,guardar centro) & Botón \\
\rowcolor{TblId}
3 & 1 & El mismo id repetido & Botón \\
\end{tabular}

\vspace{0.4em}
\chip{TblSync}{Sincronía}
\chip{TblId}{Botón}
\end{center}

\subsection{Dispositivo $\rightarrow$ teléfono}

\paragraph{\texttt{CC 11} (8 bytes)}

\begin{center}
\renewcommand{\arraystretch}{1.25}
\setlength{\tabcolsep}{4pt}
\begin{tabular}{>{\centering\arraybackslash}p{1.35cm}
                >{\centering\arraybackslash}p{1.05cm}
                >{\raggedright\arraybackslash}p{9.15cm}
                >{\centering\arraybackslash}p{2.5cm}}
\rowcolor{TblHead}
\thc{Despl.} & \thc{Tam.} & \thc{Contenido} & \thc{Rol} \\
\rowcolor{TblSync}
0--1 & 2 & Sincronía \texttt{0xCC 0x11} & Sincronía \\
\rowcolor{TblTelem}
2--3 & 2 & left, int16 LE & Telemetría \\
\rowcolor{TblTelem}
4--5 & 2 & right, int16 LE & Telemetría \\
\rowcolor{TblTelem}
6 & 1 & Flags reservados & Telemetría \\
\rowcolor{TblSum}
7 & 1 & $\sum$ bytes $2\ldots 6$ (8 bits bajos) & Checksum \\
\end{tabular}
\end{center}

\paragraph{\texttt{CC 22} (6 bytes)}

\begin{center}
\renewcommand{\arraystretch}{1.25}
\setlength{\tabcolsep}{4pt}
\begin{tabular}{>{\centering\arraybackslash}p{1.35cm}
                >{\centering\arraybackslash}p{1.05cm}
                >{\raggedright\arraybackslash}p{9.15cm}
                >{\centering\arraybackslash}p{2.5cm}}
\rowcolor{TblHead}
\thc{Despl.} & \thc{Tam.} & \thc{Contenido} & \thc{Rol} \\
\rowcolor{TblSync}
0--1 & 2 & Sincronía \texttt{0xCC 0x22} & Sincronía \\
\rowcolor{TblTelem}
2 & 1 & analog, $u8$ & Telemetría \\
\rowcolor{TblTelem}
3 & 1 & battery, $u8$ & Telemetría \\
\rowcolor{TblTelem}
4 & 1 & LED, $u8$ & Telemetría \\
\rowcolor{TblSum}
5 & 1 & $\sum$ bytes $2\ldots 4$ (8 bits bajos) & Checksum \\
\end{tabular}
\end{center}

\paragraph{\texttt{CC 33} (variable)}

\texttt{count} debe ser 4 u 8.
No envíe \texttt{CC 44}.

\begin{center}
\renewcommand{\arraystretch}{1.25}
\setlength{\tabcolsep}{4pt}
\begin{tabular}{>{\centering\arraybackslash}p{1.6cm}
                >{\centering\arraybackslash}p{1.15cm}
                >{\raggedright\arraybackslash}p{8.9cm}
                >{\centering\arraybackslash}p{2.5cm}}
\rowcolor{TblHead}
\thc{Despl.} & \thc{Tam.} & \thc{Contenido} & \thc{Rol} \\
\rowcolor{TblSync}
0--1 & 2 & Sincronía \texttt{0xCC 0x33} & Sincronía \\
\rowcolor{TblLen}
2--3 & 2 & \texttt{length} uint16 LE ($=1+\textit{count}$) & Longitud \\
\rowcolor{TblLen}
4 & 1 & \texttt{count} (4 u 8) & Cuenta \\
\rowcolor{TblChan}
5--$4{+}c$ & $c$ & $c$ muestras $u8$ (CH1\ldots CH$c$) & Muestras \\
\rowcolor{TblSum}
$5{+}c$ & 1 & $\sum$ desde offset 2 hasta la última muestra & Checksum \\
\end{tabular}

\vspace{0.35em}
{\scriptsize $c$\,=\,\texttt{count}. El checksum cubre desde el desplazamiento 2 hasta la última muestra.}
\end{center}

\section{Fuera de alcance}

\begin{itemize}[leftmargin=1.2cm]
  \item HTTP MJPEG y \texttt{/camconfig}.
  \item Asignación oficial de SSID SoftAP (despliegue, no tramas).
  \item Perfiles BLE serie que no sean NUS.
  \item Handshake opcional; el cliente no confirmará la sesión.
\end{itemize}

\section{Lista de implementación}

\begin{enumerate}[leftmargin=1.6cm]
  \item Enlazar SPP, NUS o TCP como flujo de octetos.
  \item Acumular hasta \texttt{LF}; analizar \texttt{RC:CONNECT}.
  \item Responder \texttt{RC:ACK,app,RC} o un \texttt{NAK} estructurado.
  \item Condicionar control/telemetría a \texttt{session\_ok}; fail-safe
        al desconectar.
  \item Simple: \texttt{CTRL}/\texttt{BTN}/\texttt{SET} y, si aplica,
        \texttt{DATA}/\texttt{PLOT}.
  \item Binary: \texttt{AA 55}/\texttt{BB 66} y, si aplica,
        \texttt{CC 11}/\texttt{CC 22}/\texttt{CC 33}.
  \item Documentar para qué modo de la aplicación se compiló el binario.
\end{enumerate}

\end{document}
